\subsubsection{Source identity, review and storage}

The source store preserves original bytes and an exact text derivative. The raw
document and derivative each have a SHA-256 digest. A source span records the
derivative hash, page, half-open character offsets and a hash of the quoted text.
Extraction changes create a new derivative; they do not silently move an approved
span. A reviewer can open the original page beside the extracted passage.
Scanned or poorly extracted clauses require visual confirmation before release.

The source register distinguishes circulars, annexes, FAQs, codes, bank policy
and interpretation. Dependency rows record references, definitions, amendments
and replacements; unresolved targets remain explicit. Publication date does not
automatically establish effectiveness. Applicability records identify legal
entity, regulated role, activity, product, client class and jurisdiction. A
product-provider requirement cannot enter the bank's distributor rules merely
because its title mentions a product the bank sells.

The bundle contains definitions, expressions, spans and interpretations. Its
canonical JSON hash identifies the entire immutable revision. Review decisions
live outside those bytes and approve that exact hash. The lifecycle is draft,
in review, approved, released and retired; rejection returns a revision to draft,
while an edit creates a new draft hash. Release requires distinct compliance and
engineering reviewers, resolved blocking questions and dependencies, a supported
effective interval and passing evidence for that hash. The author cannot approve
their own revision under the proposed prototype policy.

SQLite transactions enforce review revisions and idempotency. A stale review
update is rejected rather than overwriting another review. Content-addressed
blobs are written before their database references; a missing referenced blob
is an integrity error. Audit rows record accepted transitions. Hashes detect
accidental corruption but do not make a database tamper-proof against an
administrator able to rewrite both records and hashes. Enterprise identity,
signed retention and access controls belong to the later bank deployment.

\subsubsection{A source-to-decision walkthrough}

The first demonstration imports 23EC35 and Annex 1, verifies their hashes and
opens paragraph 3.1 beside the financial rule. Its two wealth definitions are
review issues until the relevant source and data mappings are accepted. An
engineer loads the supplied draft bundle and the synthetic \texttt{f01} case:
portfolio HK\$40 million and qualifying net assets zero. The trace evaluates
both inclusive comparisons and their alternative, returning \texttt{TRUE} for
\texttt{spi.financial}. It also identifies the exact source revision, input
snapshot and engine version.

The reviewer next opens \texttt{f02}, one cent below the portfolio threshold,
with the other route still false. This returns \texttt{FALSE}. In \texttt{f07},
the portfolio is unknown and net assets are HK\$70 million, so the answer is
\texttt{UNKNOWN}. In \texttt{f06}, changing only net assets to HK\$90 million
establishes the financial condition despite the still-missing portfolio. These
cases let a compliance reader inspect the interpretation before any system rule
is exported. A proposed strict comparison is challenged with \texttt{f01}.

Once dependencies and meanings are reviewed, the evidence report is attached to
the bundle hash, reviewers approve it and the release service exports the package.
The initial production adapter is a deterministic decision API; additional bank
configuration or code generation is a later adapter against the same contract.
A changed rule creates another hash and invalidates release approval. Historic
decisions still point to the old hash and their original evidence.

\subsubsection{Service boundaries an implementer can use}

The API imports sources, creates and retrieves bundles, records reviews, releases
approved hashes, evaluates snapshots, appends events, replays histories and
compares versions. The full method/path/payload table is supplied in the contract
file. The central call is
\begin{lstlisting}
evaluate(bundle, snapshot, rule_id, valid_at, known_at,
         mode="draft")
    -> EvaluationResult
\end{lstlisting}
Its response includes the tagged outcome, exact value if any, reason codes,
missing and blocking inputs, trace, source/bundle/input hashes and engine version.
Production calls require a released exact hash; draft calls visibly retain draft
mode. An HTTP failure is different from a well-formed request whose assessment
returns an error or conflict.

Longer solver or drafting calls use persisted jobs with explicit failed,
cancelled and timed-out states. They cannot alter approval records. Official
documents and model output remain untrusted input: fetching restricts hosts and
redirects, validates content and limits resources; parsing never executes a
document's instructions. The initial service binds to localhost and uses public
sources and synthetic facts. This makes the proposed behavior implementable
before private data access or production integration is considered.
